\documentclass[journal]{IEEEtran}
\usepackage{fontspec}
\usepackage{unicode-math}
\setmainfont{TeX Gyre Termes}
\setmathfont{Latin Modern Math}
\setmonofont{Latin Modern Mono}
\usepackage{graphicx}
\usepackage{booktabs,tabularx,array}
\usepackage{enumitem}
\setlist{itemsep=2pt,topsep=4pt,parsep=0pt}
\usepackage{listings}
\lstset{basicstyle=\ttfamily\small,columns=fullflexible,keepspaces=true,mathescape=true,breaklines=true,frame=lines,framesep=4pt,aboveskip=8pt,belowskip=8pt}
\usepackage[hyphens]{url}
\usepackage[hidelinks]{hyperref}
\usepackage{caption}
\captionsetup{font=small,labelfont=bf,labelsep=period}
\providecommand{\mdlgwhtsquare}{$\square$}
\providecommand{\tightlist}{\setlength{\itemsep}{2pt}\setlength{\parskip}{0pt}}
\renewcommand{\arraystretch}{1.15}
\begin{document}
\title{Supplementary Appendix: Pseudocode for Confluent Accumulation Cells}
\author{\IEEEauthorblockN{Aditya Raj}
\IEEEauthorblockA{School of Artificial Intelligence and Data Science\\Indian Institute of Technology Jodhpur, Jodhpur, India}}
\maketitle
\appendix

\hypertarget{pseudocode}{%
\section{Pseudocode}\label{pseudocode}}

The four operations of CAC, in the compare-and-swap form. tag and val
unpack a cell; pack builds one.

\begin{lstlisting}
deposit(dst, x):
  (class, neg, m, q) = decode(x)
  # finite x = $\pm m\cdot 2^q$ grid units
  if class == zero:
    update(dst, 0, 1, neg ? $2^{28}$ : 1); return
  if class == inf:
    update(dst, NB mod K, NB+2, neg ? $2^{28}$ : 1); return
  if class == nan:
    update(dst, (NB+1) mod K, NB+3, 1); return
  for each non-zero base-$2^W$ digit g
      of $m\cdot 2^q$, in bin j:
    update(dst, j mod K, j+2, neg ? -g : g)

update(dst, slot, t, v):
  old = load(cell[dst][slot])
  loop:
    if tag(old) > t: return   # Lemma 2
    new = (tag(old) == t) ? pack(t, val(old)+v)
                          : pack(t, v)
    seen = compare_and_swap(cell[dst][slot], old, new)
    if seen == old: return
    old = seen

readout(dst):
  S = 0
  for each slot whose tag t is a finite bin tag:
    S += val(cell) $\cdot\, 2^{(t-2)W}$  # exact
  if NaN seen, or both +Inf and -Inf: return NaN
  if an infinity was seen: return it
  if S == 0: return (only -0 deposited) ? -0 : +0
  return round_to_nearest_even(S)

merge(A, B):   # multi-device
  for every cell c: A[c] = A[c] $\oplus$ B[c]
\end{lstlisting}

\hypertarget{artifact-contents}{%
\section{Artifact contents}\label{artifact-contents}}

\begin{table}[!htbp]\centering\small\setlength{\tabcolsep}{4pt}
\caption{Contents of the artifact.}\label{tab:24}
\begin{tabularx}{\linewidth}{>{\raggedright\arraybackslash\hsize=1.000\hsize}X>{\raggedright\arraybackslash\hsize=1.000\hsize}X}
\toprule
Path & Content \\
\midrule
include/cac/common.h, fp32.h, bigint.h, exact.h & Random generator, IEEE decoding, big integers, rounding and the oracle \\
include/cac/accum.h & CAC and all baselines \\
include/cac/vgpu.h, runner.h, workload.h & Virtual GPU, schedulers, sharding, workloads \\
include/cac/train.h, threads.h & Training experiment and the real-thread backend \\
src/tests.cpp, src/experiments.cpp, src/main.cpp & 55 tests, experiments E1 to E8, command line \\
scripts/crosscheck.py, scripts/summarize.py & Independent re-derivation; tables and figures \\
results/ & Raw CSV files, tables.md, figures \\
\bottomrule
\end{tabularx}
\end{table}
\end{document}
